\section*{Chapter \ref{ch:g9:international-trade} --- International Trade}

\begin{solution}{exo:g9:international-trade:1}
$59.3\,\%$ in 1965 and $87.3\,\%$ in 2024.
\end{solution}

\begin{solution}{exo:g9:international-trade:2}
$450 - 520 = -70$ million dalens: a \omterm{def:g9:international-trade:trade-deficit}{trade deficit} of 70 million.
\end{solution}

\begin{solution}{exo:g9:international-trade:3}
About $56\,\%$ in 1990 and about $4.6\,\%$ in 2022.
\end{solution}

\begin{solution}{exo:g9:international-trade:4}
$150 + 30 = 180$ dalens.
\end{solution}

\begin{solution}{exo:g9:international-trade:5}
Without a \omterm{def:g9:international-trade:tariff}{tariff}: $70 - 20 = 50$ thousand metres. With a \omterm{def:g9:international-trade:tariff}{tariff} of 3 the \omterm{def:g7:money:price}{price} is 9: buyers want
$100 - 45 = 55$, weavers make $45 - 10 = 35$, \omterm{def:g6:a-connected-world:import}{imports} are $55 - 35 = 20$ thousand metres.
\end{solution}

\begin{solution}{exo:g9:international-trade:6}
$3 \times 20 = 60$ thousand dalens a year.
\end{solution}

\begin{solution}{exo:g9:international-trade:7}
The home \omterm{def:g7:money:price}{price} rises from 8 to 10 dalens. Buyers want $100 - 50 = 50$, weavers make $50 - 10 = 40$: \omterm{def:g6:a-connected-world:import}{imports}
fall from 30 to 10 thousand metres. They do not halve: they fall to a third, because each extra dalen of
\omterm{def:g9:international-trade:tariff}{tariff} both lowers the \omterm{def:g7:demand-and-supply-first-look:quantity-demanded}{quantity demanded} and raises the quantity made at home.
\end{solution}

\begin{solution}{exo:g9:international-trade:8}
With Kestria: $40\,231 - 34\,739 = 5\,492$ million dalens, a surplus. In total: $129\,777 - 128\,664 =
1\,113$ million, a surplus. (So Brenland has a deficit of $5\,492 - 1\,113 = 4\,379$ million with its other
partners taken together.)
\end{solution}

\begin{solution}{exo:g9:international-trade:9}
It contradicts the assumption that the country's buyers cannot move the world \omterm{def:g7:money:price}{price}. The country is likely
a large buyer of the good: when its buyers buy less, foreign sellers lower their \omterm{def:g7:money:price}{price}, and they bear half
of the \omterm{def:g9:international-trade:tariff}{tariff}.
\end{solution}

\begin{solution}{exo:g9:international-trade:10}
Buyers: $3 \times 55 = 165$ thousand dalens. Weavers: $3 \times 35 = 105$ thousand dalens. State:
$3 \times 20 = 60$ thousand dalens. And $105 + 60 = 165$.
\end{solution}

\begin{solution}{exo:g9:international-trade:11}
In each country, imported goods become cheaper, so buyers buy more of them and the quantity traded rises.
Producers who compete with \omterm{def:g6:a-connected-world:import}{imports} in each country sell less or at lower \omterm{def:g7:money:price}{prices}; producers who \omterm{def:g6:a-connected-world:export}{export} sell
more abroad.
\end{solution}

\begin{solution}{exo:g9:international-trade:12}
The chart shows shares of the value of \omterm{def:g6:a-connected-world:export}{exports} in two years, not the reasons. For example: more schooling
and more machines in Korea; changes in what the rest of the world wanted to buy; changes in the \omterm{def:g7:money:price}{prices} of
the goods (a share of value changes when \omterm{def:g7:money:price}{prices} change, even if quantities do not); government decisions.
The chart cannot tell these apart.
\end{solution}

\begin{solution}{pb:g9:international-trade:1}
\textbf{1.} At 10 dalens: demanded $900 - 300 = 600$, supplied $200 - 100 = 100$ thousand metres. At 20:
demanded $300$, supplied $300$.
\textbf{2.} $900 - 30p = 20p - 100$ gives $1\,000 = 50p$: 20 dalens and 300 thousand metres.
\textbf{3.} $600 - 100 = 500$ thousand metres.
\textbf{4.} $10 \times 100 = 1\,000$ thousand dalens, that is 1 million dalens a year.
\textbf{5.} $10 + 4 = 14$ dalens.
\textbf{6.} Demanded $900 - 420 = 480$; supplied $280 - 100 = 180$ thousand metres.
\textbf{7.} $480 - 180 = 300$ thousand metres.
\textbf{8.} $4 \times 300 = 1\,200$ thousand dalens.
\textbf{9.} $4 \times 480 = 1\,920$ thousand dalens.
\textbf{10.} $4 \times 180 = 720$ thousand dalens.
\textbf{11.} $1\,200 + 720 = 1\,920$: the same. Every extra dalen buyers pay on the metres they still buy goes
to the state or to the weavers.
\textbf{12.} Buyers buy $600 - 480 = 120$ thousand metres fewer; weavers make $180 - 100 = 80$ thousand
metres more.
\textbf{13.} Besides the 1,920 thousand dalens, buyers lose the 120 thousand metres they would have bought at
10 dalens and no longer buy. And the weavers' 720 is not all gain: making the 80 thousand extra metres costs
them between 10 and 14 dalens a metre, when the same cloth could be bought abroad for 10. These two losses
go to nobody, so buyers lose more than the weavers and the state gain.
\end{solution}
